\section{n1. Independent current proof audit and reusable all-support canonical calculus; normal higher-dimensional CE remains unproved.}
\textit{Supports: G331.}

The sole mathematical input is ENTIRE qualified G331, with its mandatory annotation, full source ledger and exact scope. The 422 actual FINAL fields, assembled notes, distinct expanded proofs, source/scope records, programs and compiler observations are DATA. This note reconstructs the envelope arguments used below and records the independent current audit. It does not grant historical or recursive reading credit.

Write C and D for G331's independently geometric continuous uniform cone and actual decreasing-C-net completion. Its positivity is actual negative effective finite REAL Cartier presentations on directed ACTUAL primary blowups, tested on EVERY contracted proper integral curve. Its proved comparison with ideal profiles is used; no ideal order completion is newly named geometric. All prime supports, quotient fraction-field evaluations, full normalization fibres, unit/zero/bottom cases and primary gauges remain present. Raw geometric and ideal cores need not coincide. Arbitrary-model cofinality, literal radial identification and an envelope theorem are not granted by G331.

Here is a reusable canonical lemma, proved independently. On compact X=V\_a(R), for finite continuous signed f define A\_0=R and A\_m=\{g in R:v(g)+m f(v)\textgreater{}=0 for every v in X\}. Ultrametricity makes these ideals, a power of a lies in each, and A\_m A\_n is contained in A\_(m+n). If f is negative somewhere every A\_m is proper primary; if f\textgreater{}=0 it is R. An integral-dependence equation has a uniquely least-valued leading term whenever a finite v(g) violates its bound, so A\_m is integrally closed. Infinite values satisfy the tests already. The primary/unit profile h\_m=-v(A\_m)/m is continuous, finite and below f. On ACTUAL Bl\_(A\_m), its divisor -E\_m/m is negative effective and relatively ample, with positive degree on EVERY contracted integral curve before any limiting statement. The unit case is zero.

The raw supremum p(f) over C minorants equals that over D minorants and equals sup\_m h\_m. Indeed G331 gives uniform rational primary/unit approximation in C. A small negative rational gauge correction retains the obstacle inequality; common denominators turn it into ideal multiplication. Any rational competitor -v(I)/r\textless{}=f has I contained in A\_r. Conversely each h\_m is a competitor. For g=inf\_i g\_i in D with g\textless{}=f, the increasing directed open sets \{g\_i\textless{}f+epsilon\} cover X. Compactness gives one index valid throughout X; g\_i-epsilon is a C minorant of f. Therefore g\textless{}=p(f)+epsilon; let epsilon decrease. This also handles partial bottom. The constant -max(0,-min f) supplies a finite lower barrier. No positivity or USC of p is presumed.

Let M\_1=||v(A\_1)|| and H\_N=max\_(1\textless{}=r\textless{}=N)h\_r. If m=kr+s, 0\textless{}=s\textless{}r, grading gives h\_m\textgreater{}=(kr/m)h\_r-s M\_1/m\textgreater{}=h\_r-r M\_1/m. Thus H\_N-N M\_1/m\textless{}=h\_m\textless{}=p. Every degree converges pointwise to the finite lower semicontinuous p. Compact increasing strict-test covers show that p is continuous iff H\_N converges uniformly, iff h\_m converges uniformly. In that case p belongs to C and is the greatest C/D minorant. Negative gauge shifts prove ||p(f)-p(F)||\textless{}=||f-F||. Finite homogeneous generation forces continuity: if homogeneous generators g\_j have degrees d\_j, evaluation of every homogeneous monomial gives v(A\_m)/m\textgreater{}=min\_j v(g\_j)/d\_j; powers of the generators attain the infimum over m. Hence p=-v(J)/L, with L divisible by all d\_j and J=(g\_j\textasciicircum{}(L/d\_j)). If f is negative somewhere J is proper; it is unit at each nonmaximal prime because A\_1 is, so J is primary. This implication suffices for the non-FG deductions below. Non-FG alone is not a counterexample to CE.

Under a primary change b, d(v)=v(b) is bounded, positive and continuous, T(v)=v/d(v), and f\_b(Tv)=f\_a(v)/d(v). Each defining inequality divides by d(v), so canonical ideals are unchanged, and profiles and raw envelopes divide by d(v). Order bijections preserve existence/nonexistence of greatest minorants. G331's all-support completion homeomorphism and exact primary contraction transport the same lemma. A decreasing net of continuous obstacles having C envelopes has a greatest D envelope given by the actual decreasing net of their envelopes: every competitor is below each envelope. No increasing completion or arbitrary net flattening is used.

The normal-surface theorem was rejudged at its precise published domain. For an excellent NORMAL local DOMAIN of dimension TWO, normalize the blowup of a finite product of primary ideals including a. The result is a normal projective integral surface, regular off the closed fibre, with finite singular locus and normal singular local completions. These verify exactly condition (4) of Stacks 54.14.5. Its entire 54.14.3--5 proofs give a finite regular refinement by normalized blowups. The complete-normal-surface resolution, rational/double-point steps, completion transfer and gluing remain deeper published dependencies, not an independently recursive reconstruction.

Composition ampleness 29.38.7--8 combines negative tautological bundles into O(-E) ample for an effective exceptional E of full support. Normal proper birational 37.53.6 gives H0(Y,O)=R. This particular refinement is an ACTUAL primary blowup: finitely many affine nonvanishing section charts of powers of O(-E) cover Y; equalize degrees. The full 28.27.11 section-localization proof writes a finite set of affine algebra generators as t/s\_i\textasciicircum{}k. In a common higher degree use s\_i\textasciicircum{}K and t s\_i\textasciicircum{}(K-k). Their ratios generate each affine chart, giving a projective closed immersion. The sections lie in R, generate an invertible ideal J O\_Y, and J is proper primary since it is unit off the marked point. The immersion factors into integral Bl\_J and contains its dense punctured open; its defining ideal is zero in that integral scheme. This proves the needed finite-list refinement, not general model cofinality.

On regular Y the exceptional prime curves E\_i are actual effective Cartier. Put E\_a=sum a\_i E\_i, a\_i\textgreater{}0, D=sum d\_i E\_i, and Q\_ij=deg\_(E\_i)O(E\_j), over the marked residue field. Off-diagonal entries are nonnegative and symmetric by local effective intersection lengths. Every contracted integral curve is one of these curves. Feasible z satisfy z\textless{}=min(d,0), Qz\textgreater{}=0. A sufficiently negative -M a is feasible because O(-E\_a) is generated. Coordinate maxima preserve feasibility: in a given row keep one input's diagonal coordinate and increase its nonnegative off-diagonal terms. In the compact box above the barrier, coordinate maximizers and their finite maximum give a greatest z*. Every feasible vector is dominated by one in that box after joining the barrier. Therefore Z*=sum z\_i* E\_i is actual negative effective REAL Cartier and curve-nef.

For any primary competitor -v(I)/s\textless{}=f\_D, remove its exceptional divisorial part F\_I. Normal height-one intersection puts every generator of I in H0(Y,O(-F\_I)). On EACH exceptional CURVE some generator has nonzero generic restriction after removal, so its line-bundle degree is nonnegative. Thus -F\_I/s is feasible and \textless{}=Z*. For any support p, the valuation ring of Frac(R/p) has its proper centre on Y and a generic point off the fibre. Pull back the exceptional invertible lattice with that generic trivialization. Regular generator ratios give -v(I)/s\textless{}=f\_(-F\_I/s)\textless{}=f\_(Z*) even if some generators are killed. The canonical lemma proves p(f\_D)=f\_(Z*). Signed primary differences form a separating lattice, with min(vI-vJ,vK-vH)=v(IH+KJ)-v(JH); (g)+a\textasciicircum{}N separates all supports and a gives constants. Compact lattice approximation and the Lipschitz bound therefore prove CE for ALL continuous signed obstacles on these normal surfaces, over arbitrary residue fields and in equal or mixed characteristic. No finite generation or rate follows.

The independent BFJ application remains ONLY at a smooth affine CLOSED germ over an algebraically closed characteristic-zero field. For a signed model F fixed by a retraction, the max-directed nef coefficients below F, with a negative barrier, have bounded weak-nef supremum W. BFJ 3.11 supplies substantive CLOSED-simplex continuity; the extension is the infimum over ALL retractions. Divisorial bounds, continuity on simplices and refining the fixed model put H\_W below F on all supports. G331's exact BFJ D comparison bounds every competitor by H\_W. Canonical ideals are L-infinity(mW); truncating (g)+mhat\textasciicircum{}N with N\textgreater{}=m max(0,-min F) passes divisorial tests to all supports, and G331 contracts the primary ideals. For W nonzero, 3.14 gives strict negativity; fixed-complex thinness gives A\_thin\textless{}=-CW. The WHOLE 3.18 proof works for real t\textgreater{}C; use t=m+C to obtain (1+C/m)W\textless{}=W\_m\textless{}=W and consequently 0\textless{}=H\_W-h\_m\textless{}=C max(0,-min F)/m on EVERY support by closed-simplex continuity and all-retraction infima. Uniform model density and Lipschitz give smooth CE at these hypotheses. Zero uses the unit ideal; bottom is separate; DVR and empty cases are direct. Higher-dimensional positive/mixed-characteristic source CE is not inferred.

The current conductor arguments were independently checked. If the FULL reduced conductor-quotient boundary restriction is injective, G331's finite singleton stratum and infinite quotient-fibre rule make normalization boundary restriction a homeomorphism, with positive order/isometric C/D descent. Envelopes transport and CE is equivalent on the normal centres. Primary conductors suffice. Finite radicial reduced quotients suffice too: a support has its unique support prime and purely inseparable quotient field extension, forcing w(t)=v(t\textasciicircum{}(p\textasciicircum{}e))/p\textasciicircum{}e. No perfection or uniform exponent is needed. Canonical ideals contract; extension equality or inherited rates are not claimed.

For a numerical-semigroup ring k[t\textasciicircum{}H,z]\_0 with conductor exponent c, normalization is k[t,z]\_0 and conductor t\textasciicircum{}c S. A coefficient at exponent h\textless{}c cannot be conductor: multiplication by t\textasciicircum{}(c-1-h) exposes the forbidden exponent c-1. Localization preserves this finite-module conductor. Reduced quotients are the same z-ring, so transport gives CE in the BFJ smooth range in every dimension and on normal surfaces over every field. The mixed cusp over Z\_p has the same reduced quotient and normal surface source. For H=\textless{}2,3\textgreater{}, b=z\_1+t and A=w((b,t\textasciicircum{}2,z\_2,...)), the actual primary K\_m=(b\textasciicircum{}m(z\_1-mt),t\textasciicircum{}2 b\textasciicircum{}m,t\textasciicircum{}(2m+2))+(z\_2,...)\textasciicircum{}m has no t-linear term in any characteristic. With s=w(t),u=w(z\_1),B=w(b),U=min\_(j\textgreater{}=2)w(z\_j), min(2s,u,U)=1, A=min(B,2s,U), all generators have value\textgreater{}=mA. If U=A there is equality. Otherwise, 2s\textless{}=B forces s=u=1 and the last generator gives mA+2; B\textless{}2s and s\textless{}=u gives s\textless{}=1 and the second gives mA+2s; s\textgreater{}u gives B=w(z\_1-mt)=u=1 and the first gives mA+1, even when m=0 in k. Thus the all-support error is at most 2/m. In characteristic zero, valuations killing t with z\_1=1 and monomial (t,b) weights (1,2) give A=1,2. A primary generator attaining z-order e has a nonzero z\textasciicircum{}e coefficient and no t-linear term, so its degree-e part cannot be a scalar b\textasciicircum{}e; after z=b-t some t\textasciicircum{}(e-j)b\textasciicircum{}j, j\textless{}e, survives with weight\textless{}=2e-1. Different total degrees cannot cancel it. Every nonzero finite REAL positive primary combination has strict ratio\textless{}2, while A has ratio2. This proves CE without raw attainment and without FG, not CE failure.

The retained threefold example remains a diagnostic: blow up the regular origin and two rational exceptional points P,Q. For I=(z)+m\textasciicircum{}2, the x-chart is x(w,x), hence its divisorial part is E\_1+2E\_2+2E\_3. Its negative restriction on E\_1=Bl\_(P,Q)P\textasciicircum{}2 is H-e\_P-e\_Q and the joining line has degree -1; the other exceptional planes have degree1. The same ample-section/Rees argument makes this model actual primary. A generic section on an exceptional SURFACE does not imply all-curve-nefness. The genuine primary profile remains geometric on its own blowup. This is not normal CE failure.

For current 422/420 crossings, direct factoring through the killed component and coefficient-order comparison prove that all competitors on the second branch are bounded by its generic order value. For 0\textless{}a\textless{}b, the compatible obstacle -a on branch1 and -a-(b-a)/w(y) on branch2 has canonical ideal m\textasciicircum{}ceil(mb)+x m\textasciicircum{}(ceil(ma)-1), raw p=-a at finite-x branch1 and -b elsewhere, and exact uniform gap ceil(mb)/m-a. The signed model (1/2)v((y)+m\textasciicircum{}2)-(3/2)v(m) has the same half-gap filtration. Effective negative primary divisors are ample before the limit. Monomial x=r,z\_j=1 converges to a conductor point and proves non-USC, no greatest C/D minorant, and failure of positivity of USC regularization. Perturbation less than 1/4 preserves the half-gap contradiction. These checked inherited results are not presented as new progress.

For regular normal SURFACE centres glued along the SAME marked DVR with principal nonzero kernels and common quotient gauge, normal CE gives H\_i. Set a=min H\_i(o\_i). Zero dichotomy gives either all H\_i=0 if a=0, or C barriers B\_i\textless{}=H\_i with B\_i(o\_i)=a by scaling negative H\_i, using constant a on a zero branch. For N\textgreater{}=n||H\_i-B\_i||, max(B\_i,H\_i-w(x\_i)/n)=max(B\_i,H\_i-w((x\_i)+a\_i\textasciicircum{}N)/n) is C and equals a at o\_i. Actual primary truncation, cone maxima/addition and full-fibre positive descent make a target minorant reaching H\_i off the conductor. Hence raw p equals H\_i there and a at the shared point; equality of every H\_i(o\_i) is necessary and sufficient for greatest C/D existence. An unequal value contradicts USC along a finite-kernel-value dense net. This proof requires distinct normalization centres; it does not assume independent product factors when they are absent. For completeness the inherited cross-component necessity follows directly: use regular conductor c\_0 and normalization idempotents e\_1,e\_2, so j=c\_0e\_1 and ell=c\_0e\_2 are actual R-elements. At a fibre alpha,beta on these components, the normal bound gives g(beta)\textless{}=g(nu)/K. Choose L\textgreater{}nu(ell), B with D\_0=1+B(1-nu(ell)/L)\textgreater{}K. The signed obstacle (B/L)v((ell)+a\textasciicircum{}L)-(B+1) is -1 on component1, but every minorant has value\textless{}=-D\_0/K\textless{}-1 at the fibre. Actual primary a\textasciicircum{}((B+1)n)+(j)a\textasciicircum{}n has profile reaching -1 on the dense finite-j first-component locus and equal -(B+1) on the other components. USC at alpha is impossible for a greatest minorant. If dim R/(p\_1+p\_2)\textgreater{}0, a valuation of this quotient lifts into the two minimal normalization factors, giving that fibre. Thus for excellent reduced local dimension\textless{}=2 with NORMAL component rings, universal CE iff every pairwise intersection is zero-dimensional: sufficiency is uniqueness of the component at every support below m and normal-surface/DVR transport. This classification does not include irreducible nonnormal rings.

Finally, on reduced curve links the normalization centres are DVRs and conductor values are finite. There are no conductor identifications on the finite branch set. Positive descent gives all nonpositive functions and envelope min(f,0). Literal classical/CLD logarithmic radial extension is a different class, retaining its mixed-zero obstruction. Nilpotents are invisible and dimension-zero reduction has empty boundary. Products, energy and full finite operators are separate targets.

Evidence and audit boundary: all five actual settled FINAL notes, every one of their five fields, and all five distinct expanded proofs were physically read in untruncated bounded outputs. Their parsed actual fields were checked against the settled FINAL file; the assembled-note version is distinct and is retained separately. The unique current preamble/abstract, complete current scope records, ledger, verdict, unchanged residual and entire 420 proof were read. The six complete surface carriers, two complete SH units and six complete BFJ carriers were read, and both SH formula images were viewed. Broad truncated binding/metadata displays and one oversize combined scope display have no whole-reading credit; the relevant scope was repaired in bounded output. G331's immediate applied sections and every mandatory annotation/ledger/scope remain attached. No recursive source or entire-book reading is claimed. The three supplied native programs were read only as inert DATA. Artifact n1.md records this argument; all exact current bodies remain in proof-companion.md and retained/. Rerun: python assemble.py; python verify\_delivery.py. Execution checks retention and finite evidence, not mathematics.

\[
 A_m(f)=\{g:v(g)+mf(v)\ge0\ \forall v\in V_{\mathfrak a}\},\quad
 h_m=-v(A_m)/m,\quad p(f)=\sup_m h_m.
\]
\[
 H_N-N\|v(A_1)\|_\infty/m\le h_m\le p(f),\qquad
 p(f)\text{ continuous}\Longleftrightarrow h_m\to p(f)\text{ uniformly}.
\]

